 \documentclass{article}
\usepackage{amsmath}
\usepackage{graphicx}
\graphicspath{ {images/} }
\usepackage[spanish]{babel}
\usepackage{bbm}
\usepackage{amsthm}
\usepackage{authblk}
\usepackage{hyperref}
\usepackage[utf8]{inputenc}
\usepackage{mathrsfs}
\usepackage{amsfonts}
\usepackage{amssymb}
\usepackage{enumerate}
\usepackage[left=3cm,right=2cm,top=2cm,bottom=2cm]{geometry}
\usepackage{epsfig}
 \renewcommand{\baselinestretch}{0.93}
 \thispagestyle{empty}
 \pagestyle{empty}
\setlength{\textheight}{53pc} \setlength{\textwidth} {36pc}
\setlength{\topmargin}{-4mm}
\usepackage{multicol}
\newcommand*{\email}[1]{%
    \normalsize\href{mailto:#1}{#1}\par
    }
\setlength{\columnsep}{1cm}
\newcommand{\N}{\mathbb{N}}
\newcommand{\calM}{\cal{M}}
\newcommand{\calP}{\cal{P}}
\newcommand{\F}{\mathbb{F}}
\newcommand{\R}{\mathbb{R}}
\newcommand{\Z}{\mathbb{Z}}
\newcommand{\Q}{\mathbb{Q}}
\newcommand{\C}{\mathbb{C}}
\newcommand{\bbH}{\mathbb{H}}


\theoremstyle{definition}
\newtheorem{ejer}{Ejercicio}
\author{Marcos Morales}
\affil{Universidad Finis Terrae}
\title{Semana 4\\}



	


\begin{document}


\begin{picture}(400, 75)
   \put(-40,15){\includegraphics[width=5cm]{UFT.png}}

\end{picture}


\begin{center}
\huge{Clase 4}
\end{center}

\begin{center}
\Large{ Marcos Morales, Camila Muñoz}
\end{center}

\begin{center}
	\large{Ecuaciones Diferenciales Ordinarias (EDO)}
\end{center}
\begin{center}
\setlength{\parindent}{0cm}\large{Clases centradas para capacitar a los estudiantes de ingeniería en Ecuaciones diferenciales ordinarias. \\}
\end{center}


\vspace{1cm}

\begin{center}
	\large{Contenidos:}
\end{center}

\setlength{\parindent}{2.95cm}-  Sustitución lineal	\\
\phantom{abefwfewefwfefffwi} -  Ecuación de Bernoullí  	\\

\vspace{6cm}


\begin{center}
\today
\end{center}
\begin{center}
Santiago, Chile
\end{center}


\pagestyle{empty}


\setcounter{page}{1}
\pagestyle{plain}
\setlength{\parindent}{0cm}





\newpage

\section{Sustitución lineal}


Dada una ecuación diferencial de la forma $y'= f(ax+by+c)$, entonces haciendo el cambio de variable $z=ax+by+c$, obtenemos una ecuación de variables separables. Notemos que se tiene $\displaystyle \frac{dy}{dx} = \frac{1}{b} \left( \frac{dz}{dx} -a \right) $\\


\textbf{Ejercicio: } Encontrar la solución de la ecuación diferencial $\displaystyle \frac{dy}{dx} =3y-x$ usando sustitución lineal.\\

\textit{Respuesta: } Notemos que esta es una ecuación lineal y por lo tanto podríamos resolverla usando la fórmula general, pero en este caso se exige que se use sustitución lineal y por lo tanto vamos a usar ese método. Hacemos el reemplazo $z = 3y-x$, entonces $\displaystyle \frac{dy}{dx} = \frac{1}{3}\left(  \frac{dz}{dx} +1 \right)  $. Luego

\begin{center}
$\displaystyle \frac{1}{3} \left( \frac{dz}{dx} +1\right) =z$
\end{center}


\begin{center}
$\displaystyle \frac{dz}{dx} =3z-1$
\end{center}

\begin{center}
$\displaystyle \frac{dz}{3z-1} =dx$
\end{center}

\begin{center}
$\displaystyle \int \frac{dz}{3z-1} = \int dx$
\end{center}

\begin{center}
$\displaystyle \frac{1}{3} \ln(3z-1) = x+c$
\end{center}


\begin{center}
$\displaystyle 3z-1 = ce^{3x}$
\end{center}

\begin{center}
$\displaystyle z= \frac{ce^{3x}+1}{3}$
\end{center}


\begin{center}
$\displaystyle 3y-x= \frac{ce^{3x}+1}{3}$
\end{center}


\textbf{Actividad 1: } Utilice la fórmula general de la ecuación lineal para encontrar la solución de la EDO anterior. Explique por qué las soluciones se ven distintas. \\ \\




\textbf{Actividad 2: } Encontrar la solución de

\begin{align*}
    \frac{dy}{dx} = \cos(x+y)
\end{align*}







\newpage

\section{Ecuaciones de Bernoulli}



\textbf{Definición: }Una ecuación diferencial es de Bernoulli de grado $n$ si es de la forma $y'+a(x)y=b(x)y^n$. Como podemos ver, se parece a una lineal salvo por el hecho de que $b(x)$ está multiplicado por $y^n$.\\

\textit{Para resolver una ecuación diferencial de Bernoulli, se aplica el cambio de variable $z=y^{1-n}$ y obtenemos una ED lineal}. \\

\textbf{Ejercicio: } Encontrar la solución de la ecuación diferencial $\displaystyle y'-5y = -\frac{5}{2} x y^3$.\\

\textit{Respuesta: } Tenemos una ecuación de Bernoulli de grado 3, por lo tanto tenemos que hacer el cambio de variable $z=y^{1-3}=y^{-2}$ y entonces $\displaystyle \frac{dz}{dx}=  -2 y^{-3}\frac{dy}{dx} $, luego  $\displaystyle \frac{dy}{dx}=  -\frac{y^3}{2}\frac{dz}{dx} $. Reemplazando esto último, obtenemos \\



\begin{center}
$\displaystyle -\frac{y^3}{2}\frac{dz}{dx}-5y = -\frac{5}{2} x y^3$
\end{center}

Dividiendo por $y^3$


\begin{center}
$\displaystyle -\frac{1}{2}\frac{dz}{dx}-\frac{5}{y^2} = -\frac{5}{2} x$
\end{center}

Y recordando que $z=y^{-2}$



\begin{center}
$\displaystyle -\frac{1}{2}\frac{dz}{dx}-5z = -\frac{5}{2} x$
\end{center}


\begin{center}
$\displaystyle \frac{dz}{dx}+ 10z =  5x$
\end{center}


Esta es un ED lineal con $a(x)=10$ y $b(x)=5x$, la solución viene dada por




\begin{center}
$\displaystyle z =  e^{-\int a(x)dx} \left( \int b(x)e^{\int a(x)dx} dx  +c\right) $
\end{center}



\begin{center}
$\displaystyle z =  e^{-\int 10 dx} \left( \int 5xe^{\int 10dx} dx  +c\right) $
\end{center}



\begin{center}
$\displaystyle z =  e^{-10x} \left( 5\int xe^{10x} dx  +c\right) $
\end{center}

Falta resolver  $\displaystyle \int xe^{10x} dx$, tomamos $u=x$ y $dv=e^{10x}dx$, entonces $du=dx$ y $v= \displaystyle \int e^{10x}dx = \frac{1}{10} e^{10x}$, luego


\begin{center}
$\displaystyle \int xe^{10x} dx = \frac{xe^{10x}}{10} - \int \frac{1}{10}e^{10x}dx$
\end{center}


\begin{center}
$\displaystyle \int xe^{10x} dx = \frac{xe^{10x}}{10} - \frac{1}{100}e^{10x}$
\end{center}

por lo tanto la soulción es

\begin{center}
$\displaystyle z =  e^{-10x} \left( \frac{xe^{10x}}{2} - \frac{1}{20}e^{10x}   +c\right) $
\end{center}



\newpage

\begin{center}
$\displaystyle y^{-2} =  e^{-10x} \left( \frac{xe^{10x}}{2} - \frac{1}{20}e^{10x} +c\right) $
\end{center}

Esta última ecuación es la solución de la ED inicial. Si queremos una solución explícita sería\\ \\

\begin{center}
$\displaystyle y =  \pm \frac{1}{\sqrt{ \left( \frac{x}{2} - \frac{1}{20} +ce^{-10x}\right)}} $
\end{center}



\textbf{Actividad: } Encontrar la solución explícita de

\begin{align*}
    y'-2y =\frac{3x}{y^2}
\end{align*}















\end{document} 